\scanpage{66}
% loc: 1–8, top
d'où l'iso $\varphi^e$, et si on suppose $L$ trivial i.e. $L^*$ non vide, \deleted{on a} (\unsure{donc|bien|idem} itou pour $D$, indépendamment \deleted{de} de l'hyp. qui l'avait fixée comme un $ku$) alors dans $D^*$ il y a une classe d'éléments modulo $k^{*2}$, privilégiée : celle formée par les $\varphi^e(x\otimes x)$, pour $x\in L^*$. Appelons-la $D^*(e)$ : \deleted{On aura $D^*(\lambda e)=\lambda D^*(e)$.}

% loc: 9
\[
 D^*(e)=\{\varphi_L^e(x\otimes x)\mid x\in L^*\}\qquad \text{pour }D=D_L
\]

% loc: 10–11
\deleted{donc}
\[
 D^*(\lambda e)=\lambda D^*(e)\qquad
 \left\{\begin{aligned}
 &\text{et bien sûr }\lambda^2D^*(e)=D^*(e)\\
 &\text{puisque }D^*(e)=k^{\times2}.u\ \text{pour}\\
 &u\in D^*(e)
 \end{aligned}\right.
\]

% loc: 12–14
\deleted{La subdivision de $D^*$ en des \unsure{sous}} On peut dire que pr le $k^\times$-torseur $D^*$, on a un iso can.
\[
 D^*/k^{\times2}\simeq\underline{\omega}^*/k^{*2}
\]

% loc: 15–20
\deleted{\illegible{1}} Supposons que $ {}_2\operatorname{Pic}(k)=0$, donc pour \unsure{tout|toute} \deleted{droites $L$ de} $u$ dans le cône $C^*$ des éléments de $M$ non dégénérés de carré nul, (i.e tels que $\operatorname{Tr}u=\det u=0$) la droite $L_u$ correspondante est non vide. Alors On définit \deleted{$C^*$}
\[
 C^*(e)=\bigcup_D D^*(e)
\]

% loc: 21–23, bottom
et on a ceci :
\[
 \left\{\begin{aligned}
 C^*(\lambda e)&=\lambda C^*(e)\\
 \lambda^2 C^*(e)&=C^{\unsure{2|*}}(e)
 \end{aligned}\right.\quad
 \begin{aligned}
 &\text{donc }C^*(e)\ \text{ne dépend \unsure{que|pas}}\\
 &\text{de }e\bmod k^{\times2}\deleted{\unsure{k^*}}
 \end{aligned}
\]
% archive sheet number: 65

